\textbf{Result.} The contact rule can be tied to the theory's own response cell through three steps. Each is a stated premise, and each is backed by an identity proved in the Lean library.
\begin{enumerate}
\item \textbf{The occupation is the stock's excess field in baryon units.} Section~15 gives the excess flux $J=n_B(x)\,g_b$, so each quantum of the cell carries the flux $g_b$. If the stock carries the response (Appendix O.3), the cell occupation is $n=g_s/g_b$.
\item \textbf{The free energy is the Bose relative entropy.} Its weight per volume is fixed by equipartition. The excess free energy of a cell is $\Theta D(n\Vert n_*)$, the relative entropy of geometric distributions; it is non-negative and vanishes only at $n_*$. Its Gaussian limit equals the quadratic field mismatch $(\delta g)^2/8\pi G$ exactly when the cell free-energy density is $\mathcal N\Theta=g_b^2n_*(1+n_*)/4\pi G$. In that case the cells' static susceptibility, already the $\chi_0=n_*(1+n_*)/\Theta$ of the manuscript's two-pole response, equals the gravitational compliance of the excess field. The quadratic functional used in the simulations is then the theory's own functional to leading order, with no free coefficient.
\item \textbf{Greedy descent is the zero-temperature limit of local detailed balance.} Orbit-population transitions in contact with the horizon bath have uphill-to-downhill rate ratio $\exp(-\Delta F/\Theta)$. For any macroscopic move, $\Delta F/\Theta=\int\mathcal N\,\Delta D$ is astronomically large: the cells are about 13--60\,fm across, roughly $10^{104}$ in a galaxy volume. Uphill transitions are therefore suppressed exactly, and accept-if-$F$-decreases is the correct limit. Relaxation of the cells' own free energy releases heat to the bath, so the register pays only the mechanical $\Delta K$, as implemented.
\end{enumerate}
Replacing the quadratic mismatch by the full nonlinear Bose functional in the 3-D simulations gives the same outcome. The Milky Way model sits 0.019--0.024\,dex rms from the RAR, the dwarf 0.025--0.030\,dex, and the budget-limited massive spiral 0.051\,dex; contacts conserve momentum and angular momentum to roundoff. Twenty-eight statements, thirteen ported from the bridge note and fifteen new, are in \texttt{Entropic/ContactSelection.lean}, with standard axioms only.

Three things remain open. The rate prefactor is free. The cluster gate is not derived: the theory's candidate is that cells exist only on active carriers. The back-reaction of the cell free energy on baryons away from equilibrium still needs a common-action audit.

\section{The occupation}
The cell has gap $\Theta x$ with $x=\sqrt{g_b/a_0}$, mean occupation $n_*=1/(e^x-1)$, and excess flux $\mathbf J=n_B(\widetilde x)(\hat{\mathbf e}\cdot\nabla\Phi_b)\hat{\mathbf e}$ (Section~15). At fixed baryons, $\partial J/\partial n=g_b$: the occupation has no observable effect other than the excess field it sources, at one baryonic field unit per quantum.

Appendix O.3 makes the stock the carrier of that response. The stock field in baryon units is then the only available measure of the occupation:
\[
n\equiv\frac{g_s}{g_b}\quad\text{(spherical)},\qquad
n\equiv\frac{\hat{\mathbf e}\cdot\nabla\Phi_s}{\hat{\mathbf e}\cdot\nabla\Phi_b}\quad\text{(along the director)}.
\]
This is a premise, since the manuscript does not state that the stock field \emph{is} the occupation rather than being sourced by it. It is, however, the identification forced by asking what $n$ means for the field. In nonspherical systems the curl-free part of $\mathbf J$ is what a stock field can match. The director form is therefore an identification for the projected component.

\section{The free energy and its weight}
For a geometric occupation of mean $n$, the excess free energy at fixed gap is $\Theta D(n;x)$, with
\[
D=xn-\big[(1+n)\ln(1+n)-n\ln n\big]-\ln(1-e^{-x})
=n\ln\frac{n}{n_*}-(1+n)\ln\frac{1+n}{1+n_*}\equiv\mathrm{KL}(n\Vert n_*).
\]
Write $\phi(t)=t\ln t$. Then $\mathrm{KL}=n_*\phi(a)-(1+n_*)\phi(b)$, with $a=n/n_*$ and $b=(1+n)/(1+n_*)$. Since $b$ is the convex combination $\tfrac1{1+n_*}\cdot1+\tfrac{n_*}{1+n_*}a$, convexity of $\phi$ gives $D\ge0$, with equality only at $n=n_*$. The curvature is $\partial_n^2D=1/[n(1+n)]$, the inverse of the occupation variance.

The functional on the stock is $F_D=\int\mathcal N\Theta\,D(n\Vert n_*)\,d^3x$. For $n=n_*+\delta g/g_b$ its density is $\mathcal N\Theta(\delta g)^2/[2g_b^2n_*(1+n_*)]+O(\delta g^3)$. This equals the field-mismatch density $(\delta g)^2/8\pi G$ if and only if
\begin{equation}
\mathcal N\Theta=\frac{g_b^2\,n_*(1+n_*)}{4\pi G}.
\end{equation}
Read as a statement about compliances, this says the cells' static susceptibility $\chi_0=n_*(1+n_*)/\Theta$ (per cell, in occupation units), converted to field units, equals the compliance $4\pi G$ of the gravitational field. Equivalently, the occupation fluctuations of each cell hold $\Theta/2$ of excess-field energy. This is the fluctuation--dissipation consistency between the cells and the field they source.

With $\Theta=\hbar H_0/2\pi=3.7\times10^{-53}$\,J, the condition fixes $\mathcal N$:
\begin{center}\small
\begin{tabular}{lccc}\toprule
$g_b/a_0$ & 1 & 0.1 & 0.01\\\midrule
cells per m$^3$ & $4.5\times10^{41}$ & $4.9\times10^{40}$ & $4.9\times10^{39}$\\
cell size & 13\,fm & 27\,fm & 59\,fm\\\bottomrule
\end{tabular}\end{center}
In the deep regime $\mathcal N\to\epsilon c\,g_b/(2G\hbar)$. The cell count is a consequence to be checked against the microscopic GFT, not evidence. It does show that the large-deviation limit of the cell statistics, in which $D$ is the exact rate function, is reached with enormous margin.

The convex dual of $\mathcal Q_\perp$ in Section~15 is another cost with the same minimum but a different curvature, $4x^3/[n(1+n)]$. The equipartition condition selects the Bose relative entropy, and with it the field-mismatch form near equilibrium. Away from equilibrium the two functionals differ. The nonlinear Bose form has the derivative
\[
\frac{d\psi_D}{dr}=g_b\,n_*(1+n_*)\Big[\ln\Big(1+\frac1n\Big)-x\Big],
\]
which reduces to $d\psi/dr=J-g_s$ at linear order.

\section{The contact rule}
In the slow-contact, phase-mixed regime the stock is described by orbit populations $w_i$. For a fixed orbit basis, $\partial F/\partial w_i=\langle\psi\rangle_i$, the orbit-averaged functional derivative (bridge note). The stock elements are classical, so relabelling an element's orbit changes no internal entropy. Local detailed balance with the horizon bath then requires
\[
\frac{k_{i\to j}}{k_{j\to i}}=\exp\!\Big(-\frac{F_j-F_i}{\Theta}\Big),
\qquad \frac{|\Delta F|}{\Theta}=\Big|\int\mathcal N\,\Delta D\,d^3x\Big|\gg1 .
\]
Any change in one element's orbit alters $D$ over a region containing $\sim10^{90}$ or more cells. Uphill transitions are therefore suppressed by factors far beyond any physical relevance. The rule "accept if $F$ decreases", with a free downhill rate prefactor, is the $\Theta\to0^+$ limit (Lean: \texttt{uphill\_rate\_vanishes}). Any symmetric non-negative mobility then never raises $F$ in the continuum population description (Lean: \texttt{symmetric\_flow\_dissipates}). The finite-event non-monotonicity seen in the runs comes from the self-consistent potential moving every other orbit and from the quadratic remainder, as the bridge note shows. It does not come from the bias.

\textbf{Energy.} The descent lowers the cells' own free energy, and passive relaxation of a system's free energy toward its bath releases heat. The register therefore pays only the mechanical $\Delta K$ of each impulse, as implemented. The bridge note's heat cost $Q=-\Theta\lambda\Delta D$ applies to an externally imposed bias and is not needed here.

\textbf{Mechanics.} Contacts are equal and opposite impulses along the pair separation, which conserve $\mathbf P$ and $\mathbf L$ (Lean: \texttt{pair\_impulse\_conserves\_angular\_momentum}). A cooling impulse cannot remove more than the pair's radial relative kinetic energy $\mu u^2/2$ (Lean: \texttt{central\_cooling\_bound}).

\section{Simulations with the Bose functional}
The setup is \texttt{descent3.py} with $\psi_D$ replacing $\psi$: 3-D impulses, equilibrium NFW stock, 12\,Gyr, $N=4000$, gate 175\,km\,s$^{-1}$ (open for all three). RMS offset from the RAR, 8--12\,Gyr:
\begin{center}\small
\begin{tabular}{lccc}\toprule
Case & quadratic $E_{\rm mis}$ & Bose $F_D$ & $F_D$ end/start, rises/24\\\midrule
Milky Way, 1 proposal/orbit, seeds 1--2 & 0.017--0.020 & 0.019--0.024 & 0.010--0.046, 7\\
Dwarf, 3 proposals/orbit, seeds 1--2 & 0.029--0.033 & 0.025--0.030 & 0.74--0.75, 8--10\\
Massive spiral, light halo & 0.049 & 0.051 & 0.055, 2\\\bottomrule
\end{tabular}\end{center}
The outcome does not depend on using the quadratic approximation. Energy closes to $\le10^{-3}V_f^2$ in all runs, and $\mathbf P$, $\mathbf L$ are conserved to $\le5\times10^{-17}$.

\section{The cluster gate}
A frequency cutoff fails: at equal acceleration, larger systems have lower orbital frequencies and pass more of the response (bridge note). The theory's candidate is the carrier support of the cells. $\mathcal N$ is defined only on active carriers. These are retained, dynamically coherent source records, and their order melts where carriers overlap (Section~15). Where no active carrier is present, $\mathcal N=0$, there is no bias, and the stock evolves as collisionless matter. A cluster's stock then stays CDM-like for two reasons:
\begin{itemize}
\item shock-heated intracluster gas is not an active carrier;
\item member records overlap throughout the core, including the BCG region, and their order melts there.
\end{itemize}
This rule is not a universal dispersion threshold, and it makes different predictions:
\begin{itemize}
\item At fixed velocity dispersion, fossil groups, with one dominant carrier, should follow the RAR more closely than ordinary groups.
\item Isolated massive ellipticals should follow it, as isolated lensing ellipticals do.
\item The transition should track carrier multiplicity and baryon phase rather than mass alone.
\end{itemize}
It has not been simulated: that needs a member-galaxy model with overlapping records.

\section{Status}
\begin{center}\small
\begin{tabular}{p{0.27\linewidth}p{0.66\linewidth}}\toprule
Now grounded & the form of the bias (Bose relative entropy); the weight of the quadratic mismatch (equipartition); greedy acceptance (zero-temperature limit of local detailed balance); no heat cost for passive descent; momentum and angular momentum conservation\\
Premises & the stock field is the cell occupation; equipartition between cells and field; orbit populations exchange with the cells (one global rate prefactor); cells exist only on active carriers\\
Open & a microscopic exchange operator and its rate; a back-reaction audit of $F_D$ on baryons away from equilibrium (its first variation vanishes at equilibrium, but transient recoil must be computed in a common action); the cluster gate as a computed carrier-support model; 3-D baryons and cosmological assembly\\
\bottomrule\end{tabular}\end{center}

\section*{Lean}
\texttt{Entropic/ContactSelection.lean} has 28 statements, all checked in the main 4.35 library with standard axioms. Thirteen are ported from the bridge note: Bose entropy and gap derivatives, affinity, dissipative sign, orbit pairing, exact field variation, finite-step descent, central-impulse energetics, heat bookkeeping and two-pole monotonicity. Fifteen are new:
\begin{itemize}
\item \texttt{boseGap\_eq\_relEnt}, \texttt{relEnt\_nonneg}, \texttt{boseGap\_nonneg}, \texttt{boseGap\_at\_nStar}, \texttt{relEnt\_self}: the gap function is a relative entropy, non-negative, and zero at $n_*$.
\item \texttt{boseGap\_second\_derivative}: its curvature is $1/[n(1+n)]$.
\item \texttt{equipartition\_matching}: the Gaussian limit of the cell functional equals the field mismatch exactly under the equipartition condition.
\item \texttt{uphill\_rate\_vanishes}: local detailed balance as $\Theta\to0^+$ suppresses every uphill transition.
\item \texttt{symmetric\_flow\_identity}, \texttt{symmetric\_flow\_dissipates}: a finite H-theorem for orbit populations.
\item \texttt{pair\_impulse\_conserves\_momentum}, \texttt{pair\_impulse\_conserves\_angular\_momentum}: conservation by pair impulses along the separation.
\item \texttt{stiff\_coupling\_amplification}: the exact algebra of constraint 2.
\item \texttt{vp\_acceleration\_scaling}: constraint 1.
\item \texttt{clip\_renormalise\_not\_minimal}: the three-shell counterexample.
\end{itemize}

\textbf{Reproducibility.} \texttt{descent3.py} (argument 9: \texttt{D} for the Bose functional), \texttt{runD.sh}, \texttt{summarize3.py}; outputs \texttt{d3\_*\_FD.json}, \texttt{log3D\_*.txt}. Lean: \texttt{Entropic/ContactSelection.lean} in the formalization repository.
